\subsection{A higher PL--TOP discrepancy on the three-sphere}

The component calculation for spheres does not extend to an unrestricted
higher PL--TOP comparison.  In the following statement
$\mathrm{TOP}(4)=\operatorname{Homeo}(\mathbb R^4;0)$ is the group of
global origin-fixing homeomorphisms with its compact-open topology, and
\[
 Q_4=\operatorname{hofib}
       \bigl(B\mathrm O(4)\longrightarrow B\mathrm{TOP}(4)\bigr).
\]
Thus $Q_4$, also denoted $\mathrm{TOP}(4)/\mathrm O(4)$ in smoothing
theory, is a homotopy fibre.  All loop spaces and positive homotopy
groups below use the identity basepoint.

\begin{theorem}\label{thm:three-sphere-higher-top-discrepancy}
For the standard PL three-sphere, with bottom fixed pointwise and
upper boundary free, there is a based weak equivalence
\[
                  P_{\mathrm{TOP}}(S^3)\simeq\Omega Q_4.
\]
Here the TOP concordance object is the singular complex of the
compact-open homeomorphism group, and the PL object uses jointly
finite PL families on the whole parameterized product.  Their
homotopy groups satisfy
\[
 \begin{gathered}
 \pi_0P_{\mathrm{PL}}(S^3)=\pi_1P_{\mathrm{PL}}(S^3)
          =\pi_2P_{\mathrm{PL}}(S^3)=0,\\
 \pi_0P_{\mathrm{TOP}}(S^3)=\pi_1P_{\mathrm{TOP}}(S^3)=0,
 \qquad \pi_2P_{\mathrm{TOP}}(S^3)\cong\mathbb Z/2.
 \end{gathered}
\]
Consequently the natural real PL-to-TOP map fails to be a weak
equivalence for this base.
\end{theorem}

\begin{proof}
Let $D\subset\mathbb R^4$ be the closed unit ball.  We use singular
complexes of topological groups throughout, and write
\[
 A=\operatorname{Sing}\operatorname{Homeo}(D;0),\qquad
 T=\operatorname{Sing}\operatorname{Homeo}(\mathbb R^4;0).
\]
Let $G$ be the simplicial group of parameterized germs at the zero
section.  A $q$-simplex is a germ of a homeomorphism between open
neighbourhoods of $\{0\}\times\Delta^q$ in
$\mathbb R^4\times\Delta^q$, preserving the parameter and fixing
the zero section.  Representatives are identified when they agree
on a neighbourhood of the entire zero section.  In particular,
compactness of the parameter simplex supplies common positive
source and image radii for each representative and its inverse.
There are restriction homomorphisms
\[
                       r:A\longrightarrow G,
                 \qquad g:T\longrightarrow G.
\]

\medskip\noindent\textit{A relative lifting statement.}
We first prove that these are Kan fibrations, with exact agreement
on the prescribed parameter faces.  The geometric input is a
continuous local ambient-extension operator.  For concentric compact
balls $C\subset\operatorname{int}U\subset U\subset\operatorname{int}D$,
there is a neighbourhood $\mathcal V$ of the inclusion in the
embedding space $\operatorname{Emb}(U,\operatorname{int}D)$ and a
continuous map
\[
 \sigma:\mathcal V\longrightarrow
       \operatorname{Homeo}(D\;\mathrm{rel}\;\partial D)
\]
such that $\sigma(v)|_C=v|_C$ and
$\sigma(\operatorname{id})=\operatorname{id}$.  Its support may be
kept in a fixed compact subset of $\operatorname{int}D$.
This is the canonical ambient extension obtained from
\cite[Theorem~5.1, p.~70; proof of Corollary~1.2, p.~79]
{EdwardsKirby1971}.  Their deformation depends continuously on the
input embedding, fixes the inclusion, ends at an embedding equal to
the identity on $C$, and is unchanged off a compact neighbourhood
of $C$ in $U$.  Composing the input with the inverse of that endpoint
and extending by the identity gives $\sigma$.  Shrinking the input
neighbourhood keeps all supports away from $\partial D$.

Let $H=\Lambda_i^q\subset\Delta^q$ be a horn.  Suppose that a
family $a:H\to\operatorname{Homeo}(D;0)$ lifts the restriction of
a given germ simplex represented by $f_t$.  There are finitely
many horn faces; equality of their parameterized germs therefore
gives, after one common shrinking,
\[
 f_t=a_t\quad\hbox{on a fixed ball, for all }t\in H.
\]
Choose a strong deformation $d_s$ of $\Delta^q$ from a retraction
onto $H$ to its identity, fixed on $H$.  A simplex collapses onto
each of its horns, so such a deformation exists.  Set
\[
 a^0_t=a_{d_0(t)},\qquad
 k_{t,s}=f_{d_s(t)}(a^0_t)^{-1}
\]
as germs.  Uniform shrinking, using joint continuity of the maps
and inverses on the compact parameter set, represents all these
germs and their inverses on common neighbourhoods.  On a smaller
common neighbourhood they satisfy the actual equalities
\[
       k_{t,0}=\operatorname{id},\qquad
       k_{t,s}=\operatorname{id}\quad(t\in H).
\]

Choose $U$ so small that every $k_{t,s}^{-1}$ is defined on $U$
and sends it into a common source neighbourhood for all the maps
$k_{t,s'}$.  Choose $C\subset\operatorname{int}U$, and then a
still smaller ball $V$ such that
\[
                  k_{t,s}(V)\subset\operatorname{int}C
                  \quad(t\in\Delta^q,\ s\in I).
\]
These choices follow from compactness and joint continuity at the
fixed zero section.  The increments
\[
          v_{t;s,s'}=k_{t,s}k_{t,s'}^{-1}|_U
\]
are continuous in the embedding topology and equal the inclusion
when $s=s'$.  Compactness of this diagonal supplies a single
$\eta>0$ such that $|s-s'|<\eta$ puts every increment in
$\mathcal V$.  In particular all those increments take values in
$\operatorname{int}D$.

Take a finite partition $0=s_0<\cdots<s_N=1$ with mesh less
than $\eta$.  Starting with $b_{t,0}=\operatorname{id}$, define
on each subinterval
\[
 b_{t,s}=\sigma(v_{t;s,s_j})b_{t,s_j},
                  \qquad s_j\leq s\leq s_{j+1}.
\]
The formulas fit at the endpoints because $\sigma(\operatorname{id})$
is the identity.  Inductively $b_{t,s}|_V=k_{t,s}|_V$: the input
to the next extension is $k_{t,s_j}(x)\in C$ for $x\in V$.
The extensions fix $0$.  On the whole horn every increment is
the identity, so $b_{t,s}$ is the identity there for every $s$.
Consequently
\[
                       \widetilde a_t=b_{t,1}a^0_t
\]
is a disk-homeomorphism family equal to $a$ on $H$, with germ
$f_t$ on a common ball.  This is the required horn lift.  For $g$,
extend the same supported corrections by the identity outside
$D$ and compose them with the original global family.  This proves
both fibration statements without assuming that the original
family itself is close to the identity.

\medskip\noindent\textit{The two fibres.}
Let $K=\ker r$.  For $0<\varepsilon<1$, write
\[
 K_\varepsilon=
 \operatorname{Sing}\operatorname{Homeo}
              (D\;\mathrm{rel}\;D_\varepsilon),
 \qquad D_\varepsilon=\varepsilon D.
\]
The kernel $K$ is the filtered union of these simplicial groups.
Restriction to the annulus, followed by a fixed product
parameterization, identifies $K_\varepsilon$ with
$P_{\mathrm{TOP}}(S^3)$.  Extension by the identity on the inner
ball is the continuous inverse.  The inner sphere is fixed and
the outer sphere is free, exactly as required.

If $0<\varepsilon<\varepsilon_0<1$, choose increasing radial
homeomorphisms $\rho_u$ with $\rho_0=\operatorname{id}$ and
\[
 \rho_u(D_{\varepsilon_0})=D_{v(u)},\qquad
 \varepsilon\leq v(u)\leq\varepsilon_0,
 \qquad v(1)=\varepsilon.
\]
For example, linearly interpolate the increasing piecewise affine
radius functions.  Conjugation by $\rho_u$ gives a homotopy within
$K_\varepsilon$ from the inclusion of $K_{\varepsilon_0}$ to an
isomorphism onto $K_\varepsilon$.  Every finite simplicial sphere
and homotopy in the filtered union uses a single common radius.
Hence
\[
          K_{\varepsilon_0}\xrightarrow{\ \simeq\ }K,
                 \qquad K\simeq P_{\mathrm{TOP}}(S^3).
\]

The global kernel $L=\ker g$ is weakly contractible.
Every finite simplicial family $h_t$ in $L$ is represented on finitely
many parameter simplices. Their identity-germ conditions therefore
give a common fixed ball $D_\varepsilon$. The contraction
\[
 c_s(h_t)(x)=
 \begin{cases}s^{-1}h_t(sx),&s>0,\\x,&s=0\end{cases}
\]
is continuous in the compact-open topology: on any fixed compact
set it is exactly the identity for all sufficiently small $s$,
uniformly in $t$, and the same assertion holds for its inverse.
It stays in the kernel, fixing $D_\varepsilon$ for $0\leq s\leq1$,
and keeps identity parameter values fixed.  This contracts all
finite homotopy representatives and proves the claim.

Write $T_0,G_0$ for the simplicial identity components.  The
fibration $g$ restricts to $T_0\to G_0$.  Every vertex of $G_0$
is joined to the identity by finitely many edges, which lift by
the fibration property.  Its fibre is $L$, so the fibre sequence
gives a weak equivalence
\[
                            T_0\xrightarrow{\ \simeq\ }G_0.
\]
No globalization of a germ in an arbitrary other component is
needed here.

\medskip\noindent\textit{Identifying the disk group.}
Restriction to the boundary gives
$\operatorname{Homeo}(D;0)\to\operatorname{Homeo}(S^3)$,
with the radial-coning section.  Its kernel is contractible by
the relative Alexander isotopy
\[
 a_u(x)=
 \begin{cases}
      u\,a(x/u),&\|x\|\leq u,\ u>0,\\
      x,&\|x\|\geq u\text{ or }u=0.
 \end{cases}
\]
Both formulas agree on the dividing sphere.  They fix the origin,
and the displacement of the map and its inverse is at most $2u$,
which gives uniform continuity at $u=0$.  Splitting by the coned
boundary map proves that boundary restriction is a homotopy
equivalence.

The actual topological three-sphere group has the homotopy type
given by the natural inclusion
$\mathrm O(4)\to\operatorname{Homeo}(S^3)$
\cite[Appendix, statements~(3) and~(16), pp.~605--606]{Hatcher1983}.
The cited appendix explicitly includes TOP, so this is not merely
the smooth Smale theorem.  Thus
\[
      \operatorname{Sing}\mathrm O(4)\xrightarrow{\ \simeq\ }A,
 \qquad \operatorname{Sing}\mathrm{SO}(4)\xrightarrow{\ \simeq\ }A_+,
\]
where $A_+$ is the orientation-preserving identity component.
Every element of $K$ fixes a neighbourhood of the origin and
therefore preserves orientation; it is exactly the fibre of
$A_+\to G_0$.

Extend a disk homeomorphism to all of $\mathbb R^4$ by its
boundary map outside the disk:
\[
                  e(a)(r\omega)=r\,a|_{S^3}(\omega)
                         \quad(r\geq1).
\]
Together with $e(a)=a$ on $D$, this defines a continuous group
homomorphism $e:A\to T$ with $ge=r$.  On $\mathrm O(4)$ it is
the usual linear inclusion.  The preceding weak equivalence and
fibration statements identify
\[
 \begin{aligned}
 K&\simeq\operatorname{hofib}(A_+\longrightarrow T_0)\\
  &\simeq\operatorname{hofib}(\mathrm{SO}(4)\longrightarrow T_0)
   \simeq\operatorname{hofib}_1(\mathrm O(4)\longrightarrow T).
 \end{aligned}
\]
In the last step an orientation-reversing orthogonal map cannot
be joined to the identity in $T$, while $\mathrm{SO}(4)$ lies
in $T_0$.  Thus precisely these source and target components
enter the based homotopy fibre.  Using the singular-group models
for classifying spaces, looping their defining homotopy-fibre
square identifies the last space with $\Omega Q_4$.  Together
with the description of $K$, this proves the asserted TOP model.

\medskip\noindent\textit{The homotopy groups.}
Only the $4$-connectedness of the stabilization
$Q_4\to\mathrm{TOP}/\mathrm O$ is needed here: it gives
isomorphisms on homotopy groups in degrees at most three.
This is established in the first part of the proof of
\cite[Theorem~8.7A, p.~128; proof on pp.~130--131]{FreedmanQuinn1990},
whose full statement is $5$-connectedness.
The corrected account
\cite[Section~1.1, Theorem~1.3 and the following paragraph,
p.~2]{GabaiGayHartmanKrushkalPowell2026}
specifies exactly the global origin-fixing group and
classifying-space homotopy-fibre convention used above and
supplies the correction needed for the final $5$-connected step.
The low stable groups are
\[
 \pi_i(\mathrm{TOP}/\mathrm O)=
 \begin{cases}
 \mathbb Z/2,&i=3,\\
 0,&0\leq i\leq6,\ i\ne3,
 \end{cases}
\]
by \cite[Essay~V, Theorem~5.5(II), pp.~251--252]{KirbySiebenmann1977}.
The TOP assertions follow by looping.

For PL, the jointly PL sphere fibre model
\eqref{eq:sphere-concordance-fiber} gives
\[
 P_{\mathrm{PL}}(S^3)\simeq
       \operatorname{hofib}(X_3\longrightarrow X_4),
       \qquad X_j=\mathrm{PL}_j/\mathrm O_j.
\]
The stabilization maps induce isomorphisms
$\pi_iX_j\cong\pi_i(\mathrm{PL}/\mathrm O)$ for $i<j$
\cite[Section~5, formula~(3), p.~37]{BurgheleaLashof1974II};
the stable groups vanish through degree six
\cite[Annex~C, p.~318]{KirbySiebenmann1977}.
The outer terms of
\[
 \pi_3X_4\longrightarrow\pi_2P_{\mathrm{PL}}(S^3)
                          \longrightarrow\pi_2X_3
\]
are therefore zero.  The same exact sequence gives vanishing
of the fundamental group.  For components, $X_3$ is connected
and $X_4$ is simply connected, so the based homotopy fibre is
connected as well.  This proves all the stated PL calculations
and the failure of the comparison on the second homotopy group.
\end{proof}
